\section{ODAC25 metadata-only composition case study}
\label{supp:odac25}

\subsection{Source, transfer and original-split reconciliation}

This descriptive application uses the filtered ODAC25 release
\citep{sriram2026odac25} through the ColabFit conversion
\citep{colabfit2026odac25}. The original data documentation is retained with
its Hugging Face revision. The pinned ColabFit revisions are:
\begin{flushleft}
Training: \texttt{064de0077b3de4745f58c58e7a34da933c8b727e}\\
Validation: \texttt{99a348267a440fbdb45725c7e790eb454c4ad796}
\end{flushleft}
We read the columns \texttt{names}, \texttt{chemical\_formula\_hill} and
\texttt{multiplicity} from all 82 configuration shards. HTTPS byte-range
requests retrieve Parquet footers and the selected column chunks; the reader
rejects a non-206 response or a mismatched Content-Range before reading the
body. It never downloads a complete shard or the original ASE-LMDB archive.
The retained ledger records every pinned source URL, selected columns, source
file size, payload bytes, range-request count and compact-cache checksum.

The metered production payload is 63,431,987 bytes. The source shards occupy
200.5 GB in total (200,522,544,000 bytes). A recorded pilot transfers 248,514
bytes; another small failed projected-column pilot and documentation/API
requests are not fully metered. Thus \ODACTransferMB{} MB is the production
payload, not an exact whole-session network counter. The expanded compact
JSON caches occupy about 1.2 GiB locally and are not included in that transfer
figure. Those caches can be regenerated through the pinned metadata requests.
The paper builds from retained compact results and assignments without network
access or trajectory data.

The mirror contains \ODACUniqueRows{} configuration rows. Multiplicities
restore the original counts: \ODACTrain{} training configurations and
\ODACVal{} validation configurations, totalling \ODACFrames{}. Each compact
record retains its original source fold. Counts reconcile exactly with the
official filtered archives, but this is not a row-by-row verification of their
atomic structures or energy corrections. The new two-fold assignments pool
both original folds. No original test archive is listed in the documentation
used here.

\subsection{Identity rules and composition representatives}

The mirror retains trajectory names but not the official
\texttt{atoms.info.mof\_name} field. We therefore identify conservative
\emph{inferred families}, not certified original MOF identities. Adsorption and
GCMC suffixes are removed; synthetic \texttt{Gen<number>} names are grouped;
six-letter CSD parents combine polymorph, defect and functionalization variants;
other source prefixes retain the documented parsing rule. Connected components
also link shared configuration aliases, identical reduced observed formulas,
and identical known reduced bare-host formulas. These rules give
\ODACFamilies{} inferred families and \ODACGroups{} immutable components.
No component is dissolved. Shared composition does not establish shared
topology, and the grouping rule may deliberately combine distinct structures.

Bare-host formulas come from bare-MOF rows or subtraction of explicit
adsorbate stoichiometries from recognized trajectory names. The latter agrees
with 97,870 corresponding bare trajectories with no formula mismatch and
recovers formulas for 423 additional inferred families. The 92 remaining
families have no resolved bare formula; none is guessed. All optimization
features consistently use the \emph{observed full-system} composition,
including adsorbates. Known-host formulas add identity links only.

For each of \ODACFormulas{} distinct reduced observed formulas, let $x_f$
be its 118-dimensional elemental atom-fraction vector. Each component $g$
receives the unweighted mean of its distinct $x_f$ vectors. The optimizer's
bounded Euclidean kernel uses standardization fitted over these component
representatives. This aggregation is an additional representation choice:
the kernel of a mean feature vector generally differs from the mean of the
underlying kernels. It is not claimed to preserve the original configuration-
level cut exactly.

\subsection{Weighted optimization and matched controls}

This case freezes the original v0.1 numerical factor source, included under
\path{odac25_case/method_source}, separately from the primary v0.2.0.dev3
experiments. A rank-256 Nystr\"om factor supplies component rows $b_g$.
If $w_g$ is the component's configuration multiplicity, the fold sum is
$p_c=\sum_{g:z_g=c}w_gb_g$, and the two-fold surrogate cut is
$p_{\mathrm{train}}^\top p_{\mathrm{val}}$. Binary linear assignment enforces
the exact integer configuration counts. Starting from seeded random linear
costs, each proposal maximizes self-excluded within-fold scores; a proposal is
accepted only when its weighted factor cut decreases. At most fifteen
proposals are allowed. Retained solver logs distinguish a valid incumbent
from a certified linear optimum; neither establishes global optimality of the
quadratic partition objective.

The control before optimization is the adapter's exact feasible initialization.
It is constructed without factor values using random linear
costs and the same count constraint; it is not a uniform sampler of feasible
group partitions. The dense control uses the same initialization, whole groups,
counts and guarded proposal heuristic with the full component kernel. It is
competitive at this size, and no low-rank speedup is inferred from this case.
Optimization-only timings are retained with their scope; they exclude metadata
aggregation, feature construction, factorization and scoring. The earlier
80/10/10 split balanced component counts and is not used in this count-matched
comparison.

\subsection{Evaluation on recorded compositions and uncertainty}

Evaluation returns to the \ODACFormulas{} recorded formula vectors, rather
than scoring only component means. Standardization is fitted once to the
complete configuration-weighted formula pool, using the same mean and sample
variance for all assignments. Thus the scoring reference differs explicitly
from the component-level optimization standardization. Zero-variance
dimensions are omitted. With configuration counts $N$, $n_T$ and $n_V$, the
global fraction is estimated from independently sampled total-pool and
cross-fold pairs; actual self-pair expectation $1/N$ is removed from the
total-pool denominator. The ratio is normalized by
$2n_Tn_V/[N(N-1)]$. Each pair mean uses 200,000 draws. Nearest-training
similarity averages 512 configuration-weighted validation queries, each
compared with \emph{every} distinct training formula.

The reported Monte Carlo standard errors describe pair/query sampling only.
They are not confidence intervals over independent frameworks, datasets or
model initializations. Three optimizer seeds produce only
\ODACDistinctAssignments{} distinct low-rank assignments: seeds zero and two
are identical, and seed one changes nine component labels. The validation
fold contains \ODACValidationGroups{} components, with the two largest
accounting for \ODACTopTwo\% of its configuration weight.

\begin{table}[htbp]\centering\begin{tabular}{llrrr}
\toprule
Method & Seed & Global ratio & Mean nearest & NN MC SE\\\midrule
Original & -- & 0.9692 & 0.8464 & 0.0104\\
Groups before optimization & 0 & 0.8726 & 0.5743 & 0.0109\\
Groups before optimization & 1 & 0.8318 & 0.5589 & 0.0111\\
Groups before optimization & 2 & 0.8248 & 0.5256 & 0.0119\\
Low-rank optimization & 0 & 0.2829 & 0.1937 & 0.0060\\
Low-rank optimization & 1 & 0.2833 & 0.1957 & 0.0060\\
Low-rank optimization & 2 & 0.2829 & 0.1937 & 0.0060\\
Dense-kernel control & 0 & 0.2868 & 0.1644 & 0.0051\\
Dense-kernel control & 1 & 0.2872 & 0.1579 & 0.0047\\
Dense-kernel control & 2 & 0.2866 & 0.1538 & 0.0046\\
\bottomrule\end{tabular}
\caption{ODAC25 individual observations. MC SE describes configuration-query sampling error only. It does not describe uncertainty over independent MOFs, projects, or prediction models.}\end{table}


Independent float64 SciPy nearest-neighbour calculations for 32 additional
sampled validation queries per arm agree with Torch within
\ODACNNCheck{} similarity. The largest returned low-rank cut discrepancy
against the exact component kernel is \ODACCutError\% of that cut; this
denominator differs from the total-kernel-mass normalization in the primary
fidelity study. Checksum verification covers saved inputs, producer scripts,
the frozen factor implementation and results. Small exhaustive tests separately
check exact integer assignment and agreement between weighted and explicitly
expanded configuration objectives.

\subsection{Scientific scope}

For seed zero, the total variation between mean training and validation atomic
fractions is \ODACShiftLow{}, compared with \ODACShiftOriginal{} for the
original split. Phosphorus presence changes from \ODACPhosphorusTrain\% in
training to \ODACPhosphorusVal\% in validation. Under the original split,
\ODACOriginalRepeated\% of validation configuration weight shares a reduced
observed formula with training. Its absence in the new split follows from the
grouping contract and must not be attributed to optimization alone. Original
composition overlap is also not proof that identical atomic geometries cross
folds. This case establishes a metadata-efficient composition-shift holdout,
with improvements beyond the imposed grouping constraints. It evaluates no
energies, forces, prediction errors, target-distribution shifts or prospective
project outcomes, and it is not a scalability test over millions of independent
optimizer entities.
